% Locality of a single coefficient update, for three basis families. Compile standalone.
% Perturbing c_2 (knot z_2 = -0.5) by delta changes the edge function by delta*psi_2(u),
% exactly, since phi is linear in c. The shaded area is that change; the blue band is the
% interval that layer-0 inputs actually occupy after tanh normalisation (measured from the
% trained checkpoints: u in [0.25, 1.0]).
\documentclass[border=4pt]{standalone}
\usepackage{amsmath,amssymb}
\usepackage{mathptmx}
\usepackage[scaled=0.90]{helvet}
\usepackage{pgfplots}
\pgfplotsset{compat=1.17}
\usepgfplotslibrary{fillbetween,groupplots}

\definecolor{ink}{RGB}{28,32,38}
\definecolor{primary}{RGB}{155,26,32}
\definecolor{accentsoft}{RGB}{247,228,229}
\definecolor{rule}{RGB}{209,213,219}
\definecolor{blueacc}{RGB}{30,90,168}
\definecolor{bluesoft}{RGB}{232,239,248}
\definecolor{muted}{RGB}{95,107,122}

% uniform cubic B-spline bump on knot spacing h = 0.5, centred at z
\pgfmathdeclarefunction{bsp}{2}{%
  \pgfmathparse{abs((#1-#2)/0.5) < 1 ? (2/3 - ((#1-#2)/0.5)^2 + abs((#1-#2)/0.5)^3/2)
     : (abs((#1-#2)/0.5) < 2 ? (2-abs((#1-#2)/0.5))^3/6 : 0)}}
% Gaussian RBF exactly as implemented: exp(-(u-z)^2 / (2 h^2)), h = 0.5
\pgfmathdeclarefunction{rbfg}{2}{\pgfmathparse{exp(-0.5*((#1-#2)/0.5)^2)}}

\begin{document}
\begin{tikzpicture}[font=\sffamily\small]
\begin{groupplot}[
  group style={group size=3 by 1, horizontal sep=1.25cm},
  width=6.1cm, height=4.6cm,
  xmin=-1, xmax=1, ymin=-1.15, ymax=1.75,
  xtick={-1,-0.5,0,0.5,1},
  ytick={-1,0,1},
  tick label style={font=\scriptsize},
  xlabel={normalised input $u$}, xlabel style={font=\scriptsize, yshift=3pt},
  axis line style={ink, line width=0.5pt},
  tick style={ink},
  axis lines=left,
  clip=true, set layers,
  title style={font=\sffamily\small\bfseries, yshift=-3pt},
]
% ------------------------------------------------------------ B-spline
\nextgroupplot[title={Cubic B-spline: compact support}]
\begin{pgfonlayer}{axis background}\fill[bluesoft] (axis cs:0.25,-1.15) rectangle (axis cs:1,1.75);\end{pgfonlayer}
\foreach \z in {-1,-0.5,0,0.5,1} {
  \addplot[ink!35, line width=0.5pt, domain=-1:1, samples=121] {bsp(x,\z)};
}
\addplot[name path=Bb, draw=none, domain=-1:1] {0};
\addplot[name path=Bt, draw=none, domain=-1:1, samples=161] {bsp(x,-0.5)};
\addplot[accentsoft] fill between[of=Bt and Bb];
\addplot[primary, line width=1.2pt, domain=-1:1, samples=161] {bsp(x,-0.5)};
\node[font=\scriptsize] at (axis cs:0,1.47) {max $|\Delta\phi|$ on data $=0.021\,\delta$};
% ------------------------------------------------------------ RBF
\nextgroupplot[title={Gaussian RBF: fast decay}]
\begin{pgfonlayer}{axis background}\fill[bluesoft] (axis cs:0.25,-1.15) rectangle (axis cs:1,1.75);\end{pgfonlayer}
\foreach \z in {-1,-0.5,0,0.5,1} {
  \addplot[ink!35, line width=0.5pt, domain=-1:1, samples=121] {rbfg(x,\z)};
}
\addplot[name path=Rb, draw=none, domain=-1:1] {0};
\addplot[name path=Rt, draw=none, domain=-1:1, samples=161] {rbfg(x,-0.5)};
\addplot[accentsoft] fill between[of=Rt and Rb];
\addplot[primary, line width=1.2pt, domain=-1:1, samples=161] {rbfg(x,-0.5)};
\node[font=\scriptsize] at (axis cs:0,1.47) {max $|\Delta\phi|$ on data $=0.32\,\delta$};
% ------------------------------------------------------------ Chebyshev
\nextgroupplot[title={Chebyshev: global support}]
\begin{pgfonlayer}{axis background}\fill[bluesoft] (axis cs:0.25,-1.15) rectangle (axis cs:1,1.75);\end{pgfonlayer}
\addplot[name path=CbL, draw=none, domain=-1:0, samples=2] {0};
\addplot[name path=CtL, draw=none, domain=-1:0, samples=2] {x};
\addplot[accentsoft] fill between[of=CtL and CbL];
\addplot[name path=CbR, draw=none, domain=0:1, samples=2] {0};
\addplot[name path=CtR, draw=none, domain=0:1, samples=2] {x};
\addplot[accentsoft] fill between[of=CtR and CbR];
\addplot[ink!35, line width=0.5pt, domain=-1:1, samples=121] {1};
\addplot[ink!35, line width=0.5pt, domain=-1:1, samples=121] {2*x^2-1};
\addplot[ink!35, line width=0.5pt, domain=-1:1, samples=121] {4*x^3-3*x};
\addplot[ink!35, line width=0.5pt, domain=-1:1, samples=121] {8*x^4-8*x^2+1};
\addplot[primary, line width=1.2pt, domain=-1:1, samples=161] {x};
\node[font=\scriptsize] at (axis cs:0,1.47) {max $|\Delta\phi|$ on data $=1.00\,\delta$};
\end{groupplot}

\node[font=\sffamily\scriptsize, anchor=north west, align=left, text=ink] at ([yshift=-0.55cm]group c1r1.south west)
  {\tikz{\fill[accentsoft,draw=primary] (0,0) rectangle (0.35,0.2);}\; change $\Delta\phi(u)=\delta\,\psi_2(u)$ produced by the update $c_2\leftarrow c_2+\delta$
   (knot $z_2=-0.5$; for Chebyshev the second polynomial, $T_1$)\\[2pt]
   \tikz{\fill[bluesoft] (0,0) rectangle (0.35,0.2);}\; interval actually reached by layer-1 inputs after normalisation, $u=\tanh(\text{state})\in[0.25,1]$
   \qquad
   \tikz{\draw[ink!35, line width=0.5pt] (0,0.1) -- (0.35,0.1);}\; the other four basis functions};
\end{tikzpicture}
\end{document}
